%------------------------------------------------------------------------------%
\documentclass[fleqn,utf8,aspectratio=169,12pt]{beamer}

\usepackage{integracao_e_series}

\title{Séries de Potências -- Definição}

%------------------------------------------------------------------------------%
\begin{document}

\addtitle

%------------------------------------------------------------------------------%
\section{Séries de Potências}

%------------------------------------------------------------------------------%
\begin{frame}{Séries de Potências}
  Uma série da forma
  \[
    \sum_{n=0}^\infty {\color{blue}c_n} (x-{\color{magenta}a})^n
    \pause \;=\;      {\color{blue}c_0} +
                      {\color{blue}c_1} (x-{\color{magenta}a}) +
                      {\color{blue}c_2} (x-{\color{magenta}a})^2 + \cdots +
                      {\color{blue}c_n} (x-{\color{magenta}a})^n + \cdots
  \]
  \pause
  onde ${\color{blue}c_n}$ e ${\color{magenta}a}$ são constantes e $x\in\R$ é uma variável \\[3mm]
  \pause
  é chamada de \emph{Série de Potências}
  \pause
  centrada em \({\color{magenta}a}\)

  \pause
  \vspace{\baselineskip}
  Note que o índice começa em zero
\end{frame}

%------------------------------------------------------------------------------%
\begin{frame}{Atenção}
  Séries de Potência e Séries de Taylor \emph{não} são a mesma coisa
\end{frame}

%------------------------------------------------------------------------------%
\begin{frame}{Somas Parciais}
  \onslide<+->{As \emph{Somas Parciais} de uma Série de Potências são \emph{Polinômios}}

  \vspace{0.5\baselineskip}
  \begin{align*}
    \onslide<+->{P_n(x)
                 & \;=\; \sum_{k=0}^n {\color{blue}c_k} (x-{\color{magenta}a})^k} \\[5mm]
    \onslide<+->{& \;=\; {\color{blue}c_0} +
                         {\color{blue}c_1} (x-{\color{magenta}a}) +
                         {\color{blue}c_2} (x-{\color{magenta}a})^2 + \cdots +
                         {\color{blue}c_n} (x-{\color{magenta}a})^n}
  \end{align*}
\end{frame}

%------------------------------------------------------------------------------%
\section{Exemplos}

%------------------------------------------------------------------------------%
\addtocounter{exemplo}{1}
\begin{frame}{Exemplo \theexemplo}
  Determine a soma da série
  \[
    \sum_{n=0}^\infty x^n
  \]

  \pause
  Essa é uma Série de Potências com $\;a=0\;$ e $\;c_n=1$, para $n=0, 1, 2,\dots$
  \[
    \sum_{n=0}^\infty x^n = 1 + x + x^2 + x^3 + \cdots + x^n + \cdots
  \]
\end{frame}

%------------------------------------------------------------------------------%
\begin{frame}{Exemplo \theexemplo}

  Fixando $x$, temos uma \emph{Série Geométrica}
  \pause
  com
  \[
    \alpha = 1
    \qquad
    r = x
  \]
  \pause
  converge quando \(\abs{r} = \abs{x} < 1\)

  \pause
  \vspace{\baselineskip}
  Assim no intervalo aberto $(-1,\,1)$ temos
  \[
    \sum_{n=0}^\infty x^n
    \;=\; \frac{\alpha}{1-r}
    \pause
    \;=\; \frac{1}{1-x}
  \]
\end{frame}

%------------------------------------------------------------------------------%
\begin{frame}{Exemplo \theexemplo}
  \vspace{1.5\baselineskip}
  \scalebox{0.86}{% Title: gl2ps_renderer figure
% Creator: GL2PS 1.4.2, (C) 1999-2020 C. Geuzaine
% For: Octave
% CreationDate: Tue Sep 28 14:08:49 2021
\setlength{\unitlength}{1pt}
\begin{picture}(0,0)
\includegraphics[width=17cm,height=7cm]{fig_taylor_1div1mx}
\end{picture}%
\begin{picture}(481,198)(0,0)
\fontsize{10}{0}\selectfont\put(43.3701,19.6115){\makebox(0,0)[t]{\textcolor[rgb]{0.15,0.15,0.15}{{-1}}}}
\fontsize{10}{0}\selectfont\put(428.882,19.6115){\makebox(0,0)[t]{\textcolor[rgb]{0.15,0.15,0.15}{{1}}}}
\fontsize{10}{0}\selectfont\put(231.13,62.2846){\makebox(0,0)[r]{\textcolor[rgb]{0.15,0.15,0.15}{{1}}}}
\fontsize{10}{0}\selectfont\put(231.13,97.477){\makebox(0,0)[r]{\textcolor[rgb]{0.15,0.15,0.15}{{2}}}}
\fontsize{10}{0}\selectfont\put(231.13,132.669){\makebox(0,0)[r]{\textcolor[rgb]{0.15,0.15,0.15}{{3}}}}
\fontsize{10}{0}\selectfont\put(231.13,167.862){\makebox(0,0)[r]{\textcolor[rgb]{0.15,0.15,0.15}{{4}}}}
\fontsize{10}{0}\selectfont\put(438.52,97.477){\makebox(0,0)[l]{\textcolor[rgb]{0,0,0}{{$P_1$}}}}
\fontsize{10}{0}\selectfont\put(438.52,132.669){\makebox(0,0)[l]{\textcolor[rgb]{0,0,0}{{$P_2$}}}}
\fontsize{10}{0}\selectfont\put(438.52,167.862){\makebox(0,0)[l]{\textcolor[rgb]{0,0,0}{{$P_3$}}}}
\fontsize{10}{0}\selectfont\put(438.52,196.016){\makebox(0,0)[l]{\textcolor[rgb]{0,0,0}{{$P_4$}}}}
\fontsize{10}{0}\selectfont\put(376.838,185.458){\makebox(0,0)[r]{\textcolor[rgb]{0,0,0}{{\(\dfrac{1}{1-x}\)}}}}
\end{picture}
}
\end{frame}

%------------------------------------------------------------------------------%
\addtocounter{exemplo}{1}
\begin{frame}{Exemplo \theexemplo}
  Calcule a soma da série
  \[
    \sum_{n=0}^\infty \,{\color{blue}\beta}^n\,(x-{\color{magenta}a})^n
  \]
  \(\color{magenta}a\) e \({\color{blue}\beta}\neq 0\) números reais

  \pause
  \vspace{\baselineskip}
  Série de Potências centrada em \(a\) com coeficientes \(c_n=\beta^n\)
\end{frame}

%------------------------------------------------------------------------------%
\begin{frame}{Exemplo \theexemplo}
  \[
    \onslide<+->{\sum_{n=0}^\infty \beta^n(x-a)^n}
    \onslide<+->{\;=\;\sum_{n=0}^\infty \big(\beta(x-a)\big)^n}
  \]

  \onslide<+->{
  Para cada \(x\) fixo, é uma Série Geométrica com \(\alpha=1\) e \(r=\beta(x-a)\)}

%  \vspace{\baselineskip}
  \onslide<+->{Converge quando \(\abs{r}<1\)}
  \begin{columns}
    \column{0.3\linewidth}
    \begin{align*}
      \onslide<+->{\abs{\beta(x-a)}                   & \;<\; 1                         \\[2mm]}
      \onslide<+->{\abs{(x-a)}                        & \;<\; \frac{1}{\abs{\beta}}     }
    \end{align*}
    \column{0.5\linewidth}
    \begin{align*}
      \onslide<+->{-\frac{1}{\abs{\beta}}  \;<\;  & x-a \;<\; \frac{1}{\abs{\beta}}     \\[2mm]}
      \onslide<+->{a-\frac{1}{\abs{\beta}} \;<\;  & x   \;<\; a + \frac{1}{\abs{\beta}}}
    \end{align*}
  \end{columns}
\end{frame}

%------------------------------------------------------------------------------%
\begin{frame}{Exemplo \theexemplo}
  A série converge no intervalo
  \[
    I = \left(a-\frac{1}{\abs{\beta}} \;,\; a + \frac{1}{\abs{\beta}}\right)
  \]

  \pause
  Sua soma é
  \[
    \sum_{n=0}^\infty \beta^n(x-a)^n
    \;=\;
    \frac{\alpha}{1-r}
    \pause
    \;=\;
    \frac{1}{1-\beta(x-a)}
  \]
\end{frame}

%------------------------------------------------------------------------------%
\addtocounter{exemplo}{1}
\begin{frame}{Exemplo \theexemplo}
  Calcule a soma da série
  \[
    \sum_{n=0}^\infty \left(\frac{x-2}{2}\right)^{\!\!n}
    \pause
    \;=\;
    \sum_{n=0}^\infty \frac{1}{2^n}\left(x-1\right)^n
  \]

  \pause
  Série de Potências com \(a=1\) e \(c_n=\frac{1}{2^n}\)

  \pause
  Caso particular do exemplo anterior
  com \(\beta=\frac{1}{2}\) e \(a=1\)

  \pause
  Série Geométrica com \(\alpha=1\) e \(r=\frac{x-2}{2}\)

\end{frame}

%------------------------------------------------------------------------------%
\begin{frame}{Exemplo \theexemplo}
  \onslide<+->{Série Geométrica com \(\alpha=1\) e \(r=\frac{x-2}{2}\)}
  \begin{align*}
    \onslide<+->{\abs{r} & < 1             \\[2mm]}
    \onslide<+->{\abs{\frac{x-2}{2}} & < 1 \\[2mm]}
    \onslide<+->{-1 < \frac{x-2}{2} & < 1  \\[2mm]}
    \onslide<+->{-2 < x-2 & < 2            \\[2mm]}
    \onslide<+->{ 0 < x & < 4}
  \end{align*}
  \onslide<+->{A série converge no intervalo \((0, 4)\)}
\end{frame}

%------------------------------------------------------------------------------%
\begin{frame}{Exemplo \theexemplo}
  Para \(x\in (0, 4)\)
  \[
    \sum_{n=0}^\infty \left(\frac{x-2}{2}\right)^{\!\!n}
    \pause \;=\;
    \frac{\alpha}{1-r}
    \pause \;=\;
    \frac{1}{\;1-\dfrac{x-2}{2}\;}
    \pause \;=\;
    \frac{2}{\,2-(x-2)\,}
    \pause \;=\;
    \frac{2}{4-x}
  \]
\end{frame}

%------------------------------------------------------------------------------%
%------------------------------------------------------------------------------%
\section{Lista Mínima}

%------------------------------------------------------------------------------%
\begin{frame}{Lista Mínima}
  Estudar a Seção 7.1 da Apostila

  Revisar Séries Geométricas, Seção 6.1 da Apostila

  \vfill
  Atenção: A prova é baseada no livro, não nas apresentações
\end{frame}

%------------------------------------------------------------------------------%
\end{document}
%------------------------------------------------------------------------------%
